\subsection{Lifts of Idempotents}

Lemma 1. --- Let a be an element of a ring A such that \(a - a^{2}\) is nilpotent. There exists a polynomial P in \(\mathrm{X} + (\mathrm{X} - \mathrm{X}^{2})\mathbf{Z}[\mathrm{X}]\) such that \(\mathrm{P}(a)\) is idempotent in A.

Let n be a strictly positive integer such that \((a - a^{2})^{n} = 0\) . Set \(\mathrm{P}(\mathrm{X}) = 1 - (1 - \mathrm{X}^{n})^{n}\) . The polynomial \(\mathrm{P}(\mathrm{X})\) is a multiple of \(X^{n}\) , and the polynomial \(1 - \mathrm{P}(\mathrm{X})\) is a multiple of \((1 - \mathrm{X})^{n}\) , so \(\mathrm{P}(\mathrm{X}) - \mathrm{P}(\mathrm{X})^{2}\) is a multiple of \((\mathrm{X} - \mathrm{X}^{2})^{n}\) , and we have \(\mathrm{P}(a) = \mathrm{P}(a)^{2}\) . Moreover, \(\mathrm{X} - \mathrm{P}(\mathrm{X})\) is a multiple of X and 1 - X, hence a multiple of \(X - X^{2}\) .

PROPOSITION 7. --- Let a be a two-sided nil ideal of A, and let \(\overline{e}\) be an idempotent in the ring A/a. There exists an idempotent e in A whose canonical image in A/a is equal to \(\overline{e}\) .

Let \(a\) be an arbitrary representative of \(\overline{e}\) in \(\mathbf{A}\) . The element \(a - a^2\) of \(\mathbf{A}\) is nilpotent because it belongs to \(\mathfrak{a}\) . Choose a polynomial \(\mathrm{P} \in \mathbf{Z}[\mathrm{X}]\) satisfying the conditions of Lemma 1. We have \(a - \mathrm{P}(a) \in \mathrm{A}(a - a^2)\) , so the element \(e = \mathrm{P}(a)\) of \(\mathbf{A}\) has the desired property.

\pandocbounded{\includegraphics[keepaspectratio]{images/part001_2ad691f6.jpg}}

Suppose that \(\bar{e}\) belongs to the center of the ring A/a. There does not necessarily exist an idempotent e in the center Z of A lifting \(\bar{e}\) (VIII, p.~172, Exercise 31). However, if \(\bar{e}\) belongs to the image of Z in A/a, then it lifts to an idempotent in Z because \(Z \cap a\) is a nil ideal of Z.

COROLLARY 1. --- Let M and P be A-modules and u a surjective A-linear mapping from P to M. Suppose that P is projective and that there exists a nilpotent two-sided ideal \(\mathfrak{a}\) of A such that the kernel N of \(u\) is contained in \(\mathfrak{aP}\) . Let \(\mathrm{M}'\) and \(\mathrm{M}''\) be submodules of M whose direct sum is M. Then P is the direct sum of submodules \(\mathrm{P}'\) and \(\mathrm{P}''\) such that \(u(\mathrm{P}') = \mathrm{M}'\) and \(u(\mathrm{P}''') = \mathrm{M}''\) .

Denote by B the subring of \(\operatorname{End}_{\mathrm{A}}(\mathrm{P})\) consisting of the endomorphisms f of P such that \(f(\mathrm{N}) \subset \mathrm{N}\) . Let \(\overline{B}\) be the endomorphism ring of M. For \(f \in B\) , let \(\overline{f}\) be the unique endomorphism of M such that \(\overline{f} \circ u = u \circ f\) . The mapping \(f \mapsto \overline{f}\) is a ring homomorphism from B to \(\overline{B}\) . Since the module P is projective, this homomorphism is surjective; its kernel b consists of the endomorphisms \(f \in B\) such that \(f(\mathrm{P}) \subset \mathrm{N}\) . Let n be a natural number such that \(a^{n} = 0\) . We have

\[
\mathrm{P} = \mathfrak {a} ^ {0} \mathrm{P} \supset \mathfrak {a} ^ {1} \mathrm{P} \supset \dots \supset \mathfrak {a} ^ {n - 1} \mathrm{P} \supset \mathfrak {a} ^ {n} \mathrm{P} = 0,
\]

and, for every \(f \in b\) and every integer \(j \geqslant 0\) ,

\[
f (\mathfrak {a} ^ {j} \mathrm{P}) = \mathfrak {a} ^ {j} f (\mathrm{P}) \subset \mathfrak {a} ^ {j} \mathrm{N} \subset \mathfrak {a} ^ {j + 1} \mathrm{P}
\]

because \(N \subset aP\) by assumption. We therefore have \(f(P) \subset a^{j}P\) for every natural number j and every \(f \in b^{j}\) . In particular, we have \(b^{n} = 0\) .

Let \(\varepsilon'\) be the projector of M with image M' and kernel M'\,`. By Proposition 7 applied to the ring B and the nilpotent two-sided ideal b, there exists an idempotent \(e'\) in B such that \(\overline{e'} = \varepsilon'\) , that is, \(\varepsilon' \circ u = u \circ e'\) . Set \(e'' = 1 - e'\) , \(\varepsilon'' = \overline{e''}\) , P' = \(e'(P)\) , and P'\,' = \(e''(P)\) . Then P is the direct sum of the submodules P' and P'\,', and we have

\[
u (\mathrm{P} ^ {\prime}) = u (e ^ {\prime} (\mathrm{P})) = \varepsilon^ {\prime} (u (\mathrm{P})) = \varepsilon^ {\prime} (\mathrm{M}) = \mathrm{M} ^ {\prime};
\]

the equality \(u(\mathrm{P}'') = \mathrm{M}''\) is proved analogously.

COROLLARY 2. --- Let A be a ring, and let a be a nilpotent two-sided ideal of A. If P is a projective A-module, then P/aP is a projective module over A/a, and the A-module P is indecomposable if and only if the A/a-module P/aP is indecomposable.

Let P be a projective A-module, and let \(\overline{P}\) be the A/a-module P/aP. The A-module P is zero if and only if \(\overline{P}\) is zero (Theorem 2 of VIII, p.~158). Now suppose \(P \neq 0\) . Since the A/a-module \(\overline{P}\) is isomorphic to \(\overline{A} \otimes_{A} P\) , it is projective (II, §5, No.~1, p.~279, Corollary of Proposition 4). If P is indecomposable, then so is \(\overline{P}\) by Corollary 1. Conversely, suppose that P is decomposable and nonzero, and let \(P'\) and \(P''\) be two nonzero submodules of P such that \(P = P' \oplus P''\) . By Nakayama's lemma (VIII, p.~158, Theorem 2), we have \(P' + aP \neq P\) and \(P'' + aP \neq P\) ; if \(\overline{P}'\) and \(\overline{P}''\) are the canonical images of \(P'\) and \(P''\) in \(\overline{P}\) , then we have \(\overline{P}' \neq \overline{P}\) , \(\overline{P}'' \neq \overline{P}\) , and \(\overline{P} = \overline{P}' \oplus \overline{P}''\) . This proves that \(\overline{P}\) is decomposable.
% semantic-fragments: legacy-input-seam
\relax{}
